% This must be in the first 5 lines to tell arXiv to use pdfLaTeX, which is strongly recommended.
\pdfoutput=1
% In particular, the hyperref package requires pdfLaTeX in order to break URLs across lines.
\documentclass[11pt]{article}

% Remove the "review" option to generate the final version.
%\usepackage[review]{acl}
\usepackage[]{acl}

% Standard package includes
\usepackage{times}
\usepackage{latexsym}
\usepackage{booktabs}
\usepackage{multirow}
\usepackage{graphicx}
\usepackage{subcaption}
\usepackage{amsfonts}  % for mathbb
\usepackage{enumitem}
\setlist[enumerate]{itemsep=0mm}
\newtheorem{definition}{Definition}
\newtheorem{problem}{Problem}

% For proper rendering and hyphenation of words containing Latin characters (including in bib files)
\usepackage[T1]{fontenc}
% For Vietnamese characters
% \usepackage[T5]{fontenc}
% See https://www.latex-project.org/help/documentation/encguide.pdf for other character sets
% This assumes your files are encoded as UTF8
\usepackage[utf8]{inputenc}
% This is not strictly necessary, and may be commented out.
% However, it will improve the layout of the manuscript, and will typically save some space.
\usepackage{microtype}
% This is also not strictly necessary, and may be commented out.
% However, it will improve the aesthetics of text in the typewriter font.
\usepackage{inconsolata}
% If the title and author information does not fit in the area allocated, uncomment the following
%\setlength\titlebox{<dim>}
% and set <dim> to something 5cm or larger.

\title{A Large Language Model Analysis of Media and Congressional Discourse on the ESA}
% Author information can be set in various styles:
% For several authors from the same institution:
% \author{Author 1 \and ... \and Author n \\
%         Address line \\ ... \\ Address line}
% if the names do not fit well on one line use
%         Author 1 \\ {\bf Author 2} \\ ... \\ {\bf Author n} \\
% For authors from different institutions:
% \author{Author 1 \\ Address line \\  ... \\ Address line
%         \And  ... \And
%         Author n \\ Address line \\ ... \\ Address line}
% To start a seperate ``row'' of authors use \AND, as in
% \author{Author 1 \\ Address line \\  ... \\ Address line
%         \AND
%         Author 2 \\ Address line \\ ... \\ Address line \And
%         Author 3 \\ Address line \\ ... \\ Address line}

\author{AGNES NAMYALO\\
  Department of Computer Science \\
  Boise State University \\
  \texttt{agnesnamyalo@u.boisestate.edu}
  
}

\begin{document}
\maketitle

\section{Chapter 1 - Methods: Extracting ESA Narratives}
\begin{center}
  \textit{A Comparison of LLM-Based NPF Extraction and
          Structural Topic Modelling}
\end{center}

\subsection{Text corpora}
 
We assembled two text corpora. The news corpus consists of 300
\textit{Wall Street Journal} (WSJ) articles identified through the
ProQuest database using a keyword search procedure targeting ESA listing
events. Articles were retrieved for the period 1984--2025. The
congressional corpus consists of 300 documents from the Congressional
Record retrieved via the GovInfo API for the period 1995--2025,
including floor speeches, committee statements, and extensions of
remarks referencing ESA listing events. Beginning the congressional
corpus in 1995 reflects the availability of systematic digital records
through GovInfo, and is consistent with the period following the 104th
Congress when ESA oversight activity intensified \citep{Repetto2006}.\\
\\

Put a figure\\
\\
 
Both corpora were stratified by ESA policy category to ensure
representation across the range of ESA policy activities. Ten
meta-categories were constructed: ESA Regulatory/Agency Action (news:
19.3\%, congress: 19.0\%), ESA Legislative Action (news: 17.3\%,
congress: 42.7\%), ESA Species Listing (news: 17.3\%, congress: 5.7\%),
ESA Litigation (news: 11.0\%, congress: 1.0\%), ESA Habitat/Recovery
(news: 9.7\%, congress: 7.3\%), ESA Judicial Review (news: 9.3\%,
congress: 1.0\%), ESA Species Delisting (news: 4.3\%, congress: 5.0\%),
ESA Conservation Agreements (news: 1.7\%, congress: 0.7\%), ESA Funding
(news: 1.7\%, congress: 9.0\%), and Other/Non-ESA (news: 8.3\%,
congress: 8.7\%). The disparity in Species Listing representation
between corpora --- 17.3\% in news versus 5.7\% in congress --- is
itself substantively informative and consistent with the agenda-setting
hypothesis.
 

\subsection{Data Collection}

\subsection{Outcome and predictor variables}
 
The primary analysis focuses on ESA Species Listing documents, the category with the largest disparity between news and congressional representation and therefore the most direct test of the agenda-setting
hypothesis for this specific policy process.\\

The ESA encompasses a wide range of policy activities, but species listing is among the most prominent and recurring topics in ESA discourse — it is the foundational act that triggers federal protections, generates public controversy, and consistently attracts media attention from diverse stakeholders \citep{Repetto2006}. Despite this centrality, listing appears in only 5.7\% of congressional documents compared to 17.3\% of news articles in our corpus. This three-to-one disparity between media coverage and congressional production is precisely the structural gap that agenda-setting theory predicts, and it makes listing the most direct and interpretable test of the hypothesis. The primary analysis therefore focuses on ESA Species Listing.\\

The \textbf{outcome variable} is $\text{cong}_{t}$: the count of congressional ESA listing documents produced in month~$t$, aggregated from the 300-document congressional corpus. Monthly counts were
constructed by assigning each document to its date of publication and summing within calendar months over the period January~1995 to December~2025, yielding $n = 371$ monthly observations.
 
The \textbf{primary predictor} is $\text{news\_lag}$: the count of WSJ ESA listing articles published in month $t{-}1$. The one-month lag reflects the theoretical prediction of the agenda-setting hypothesis
that media coverage precedes and influences legislative attention --- and ensures that the causal direction is identified by time ordering \citep{Granger1969}.
 
The \textbf{control variable} is $\text{cong\_lag}$: the count ofcongressional ESA listing documents in month $t{-}1$, included to capture congressional own-momentum and separate it from the media
effect. Without this control, the news coefficient would absorb autocorrelation in congressional activity that is unrelated to media influence.

\section{Statistical Approach}
Selecting the appropriate model for $\text{cong}_{t}$ requires first characterising the properties of the outcome variable, then evaluating each candidate model against those properties. We follow the protocol
of making model selection decisions on structural grounds --- based on the nature of the data generating process --- rather than on post-hoc fit comparisons \citep{Zuur2010, Gelman2007}.
 
\subsection{Properties of the outcome variable}
 
Table~\ref{tab:summary} presents summary statistics for
$\text{cong}_{t}$ across the 371 monthly observations.
 
\begin{table}[H]
  \centering
  \caption{Summary statistics for the monthly count of congressional
           ESA listing documents ($\text{cong}_{t}$),
           $n = 371$ months, January~1995 -- December~2025.}
  \label{tab:summary}
  \begin{tabular}{lrl}
    \toprule
    Statistic & Value & Note \\
    \midrule
    $n$ (months)           & 371    & January 1995 -- December 2025 \\
    Months with $\text{cong}_{t} = 0$ & 359 (96.8\%) & Structural zeros \\
    Months with $\text{cong}_{t} > 0$ & 12 (3.2\%)  & Active months \\
    Maximum                & 5      & January 2025 \\
    Mean ($\bar{y}$)       & 0.046  & Near-zero baseline \\
    Variance ($\hat{\sigma}^2$) & 0.103 & Exceeds mean \\
    Var/Mean ratio         & 2.25   & $> 1$: overdispersed \\
    Skewness               & 11.27  & Extreme right skew \\
    Kurtosis               & 157.93 & Heavy tails \\
    \bottomrule
  \end{tabular}
\end{table}
 
Four properties of $\text{cong}_{t}$ are immediately relevant to model
selection. First, the outcome is a \textit{count}: it can only take
non-negative integer values. Second, it is \textit{heavily zero-inflated}:
96.8\% of observations are zero, substantially more than would be expected
from a Poisson distribution with the observed mean
\citep{Welsh2000, Mullahy1986}. Third, it is \textit{overdispersed}: the
variance-to-mean ratio of 2.25 violates the Poisson assumption that
variance equals mean \citep{Cameron1998}. Fourth, it exhibits
\textit{rare episodic spikes}: the maximum of 5 documents in January~2025
represents a value more than 100 times the overall mean, and the
distribution is correspondingly skewed ($\text{skewness} = 11.27$,
$\text{kurtosis} = 157.93$).
 
These four properties jointly eliminate Ordinary Least Squares as a
candidate model and progressively constrain the choice toward the Hurdle
model, as we now show.
 
%% ── 3.2 ─────────────────────────────────────────────────────────
\subsection{Why ordinary least squares is inappropriate}
 
OLS models the expected outcome as a linear function of predictors, minimising the sum of squared residuals:
 
\begin{equation}
  \text{cong}_{t}
    = \alpha + \beta_4 \cdot \text{news\_lag}
             + \beta_3 \cdot \text{cong\_lag}
             + \varepsilon, \quad
  \varepsilon \sim \mathcal{N}(0, \sigma^2)
  \label{eq:ols}
\end{equation}
 
Three structural requirements of OLS are violated by $\text{cong}_{t}$.
First, OLS assumes a \textit{continuous normally distributed outcome},
but $\text{cong}_{t}$ is discrete and bounded below at zero. For April
1999, where $\text{news\_lag} = 3$ and the actual outcome was
$\text{cong}_{t} = 1$, OLS predicts:
 
\begin{equation*}
  \widehat{\text{cong}}_{t}
    = 0.038 + 0.068 \times 3 + 0.033 \times 0 = 0.242
\end{equation*}
 
This prediction is categorically impossible. No data-generating process
for whole-number counts can produce 0.242 documents. Second, OLS assumes
\textit{constant variance} (homoscedasticity) in the residuals, but
count data variance typically grows with the mean, producing the
wedge-shaped residual pattern characteristic of heteroscedastic count
data \citep{Cameron1998}. Third, the fitted regression line can produce
\textit{negative predictions} for months with low news_lag values, which
are impossible for counts. OLS is therefore structurally misspecified
for this analysis and is rejected before any data examination.
 
%% ── 3.3 ─────────────────────────────────────────────────────────
\subsection{Poisson regression: the correct family, with a violated
            assumption}
 
Poisson regression belongs to the Generalised Linear Model family
\citep{Nelder1972, McCullagh1989} and is the standard baseline model
for count data. It models the expected count through a log link
function:
 
\begin{align}
  \log(\lambda_i)
    &= \alpha + \beta_4 \cdot \text{news\_lag}_i
               + \beta_3 \cdot \text{cong\_lag}_i
  \label{eq:poisson_link} \\[3pt]
  \lambda_i
    &= \exp\!\bigl(
         \alpha + \beta_4 \cdot \text{news\_lag}_i
                + \beta_3 \cdot \text{cong\_lag}_i
       \bigr)
  \label{eq:poisson_mean}
\end{align}
 
The exponential function guarantees $\lambda_i > 0$ for all covariate
values, resolving the impossibility problem of OLS. For April~1999,
Poisson gives $\hat\lambda = 0.336$, and the resulting Poisson
distribution assigns $P(\text{cong}_{t} = 1) = 24.0\%$ --- meaning the
actual observed outcome falls within a meaningful probability region of
the model's predictions.
 
Poisson regression also confirmed the direction of the agenda-setting
signal: $\hat\beta_4 = 0.714$, corresponding to a rate ratio of
$\exp(0.714) = 2.04$. Each additional WSJ ESA listing article in the
previous month is associated with a doubling of the expected
congressional count in the current month, holding prior congressional
activity constant.
 
However, the Poisson distribution requires that the variance equal the
mean --- a property known as equidispersion:
 
\begin{equation}
  \operatorname{Var}(\text{cong}_{t}) = \lambda
  \label{eq:equidispersion}
\end{equation}
 
For $\text{cong}_{t}$, the variance ($0.103$) exceeds the mean
($0.046$) by a factor of $2.25$, confirming \textit{overdispersion}
\citep{Cameron1998, Hilbe2011}. The consequence of fitting Poisson to
overdispersed data is systematic underestimation of standard errors by a
factor of approximately $\sqrt{2.25} = 1.50$, producing p-values that
are too small and inflated Type~I error rates \citep{Ver-Hoef2007}.
Standard Poisson is therefore insufficient for this data, though it
establishes the correct model family and baseline coefficient direction.
 
%% ── 3.4 ─────────────────────────────────────────────────────────
\subsection{Quasi-Poisson: correcting the standard errors}
 
Quasi-Poisson \citep{McCullagh1989} relaxes equidispersion by
introducing a dispersion parameter $\varphi$, estimated from the Pearson
residuals:
 
\begin{equation}
  \hat\varphi
    = \frac{1}{n-p}
      \sum_{i=1}^{n}
      \frac{(\text{cong}_{t_i} - \hat\lambda_i)^{2}}
           {\hat\lambda_i}
    = \frac{919.80}{368} = 2.50
  \label{eq:phi}
\end{equation}
 
where $n = 371$ and $p = 3$ ($\alpha$, $\beta_4$, $\beta_3$). The
variance is modelled as $\operatorname{Var}(\text{cong}_{t}) = \varphi
\cdot \lambda$, and all standard errors are multiplied by
$\sqrt{\hat\varphi} = \sqrt{2.50} = 1.581$. This correction preserves
the Poisson coefficient estimates while producing honest uncertainty
quantification.
 
Quasi-Poisson is a correction applied to the Poisson likelihood after
fitting, not a re-estimation under a new likelihood. The coefficients
are therefore identical to Poisson. Quasi-Poisson treats the variance as
growing linearly with the mean (equation~\ref{eq:qp_var}):
 
\begin{equation}
  \operatorname{Var}(\text{cong}_{t}) = \varphi \cdot \lambda
  \label{eq:qp_var}
\end{equation}
 
While Quasi-Poisson addresses the standard error problem, it does not
address the structural nature of the zero observations, as we discuss
in Section~\ref{sec:zeros}.
 
%% ── 3.5 ─────────────────────────────────────────────────────────
\subsection{Negative Binomial regression: quadratic overdispersion
            and fat tails}
 
The Negative Binomial (NB) distribution \citep{Hilbe2011} provides a
structurally richer model for overdispersion than Quasi-Poisson. It
introduces a dispersion parameter $\theta$ estimated through maximum
likelihood, and models the variance as a quadratic function of the mean:
 
\begin{equation}
  \operatorname{Var}(\text{cong}_{t}) = \mu + \frac{\mu^2}{\theta}
  \label{eq:nb_var}
\end{equation}
 
The $\mu^2/\theta$ term grows faster than the mean as $\mu$ increases,
producing fat tails that assign non-trivial probability to extreme
counts. When $\theta \to \infty$, equation~\eqref{eq:nb_var} collapses
to Poisson equidispersion; as $\theta \to 0$, the variance grows without
bound, accommodating arbitrarily heavy tails. For $\text{cong}_{t}$,
$\hat\theta = 0.058$ --- a small value, consistent with the extreme
January~2025 spike ($\text{cong}_{t} = 5$) being a genuine feature of
the data-generating process rather than an outlier.
 
Unlike Quasi-Poisson, NB is a full re-estimation under the NB
likelihood. Consequently, the regression coefficients differ from
Poisson: $\hat\beta_4 = 0.486$ under NB (rate ratio $= 1.63$) compared
to $0.714$ under Poisson (rate ratio $= 2.04$). The difference reflects
NB's absorption of extreme month variance into $\hat\theta$ rather than
into the predictor coefficient estimates.
 
The choice between Quasi-Poisson and Negative Binomial is a structural
one, not a numerical one. The Quasi-Poisson dispersion parameter
$\varphi$ and the NB parameter $\theta$ are not comparable quantities
--- they operate through different likelihoods on different mathematical
scales. The structural argument for NB over Quasi-Poisson in this
application is that the quadratic variance function in
equation~\eqref{eq:nb_var} better represents a data-generating process
where extreme months --- driven by concentrated political events --- can
occur with non-trivial probability that a linear variance model cannot
assign.
 
%% ── 3.6 ─────────────────────────────────────────────────────────
\subsection{The zero structure: structural versus sampling zeros}
\label{sec:zeros}
 
Both Quasi-Poisson and Negative Binomial treat the 359 zero
observations as draws from a single distribution with a low mean.
However, 96.8\% zeros in a monthly political activity series is not
simply a consequence of a small expected count. It reflects a
distinction that is critical for model specification
\citep{Mullahy1986, Welsh2000}:
 
\begin{description}[leftmargin=2em, itemsep=4pt]
  \item[Sampling zeros] arise when the expected count is low and the
    random draw from the count distribution happens to produce zero.
    These zeros are part of the count process and are appropriately
    modelled by Poisson or NB.
  \item[Structural zeros] arise when congressional engagement with ESA
    listing does not occur in a given month because the issue has not
    crossed the attention threshold --- regardless of news coverage.
    These zeros come from a separate threshold process and cannot be
    appropriately modelled by the count distribution alone.
\end{description}
 
The distinction is theoretically motivated. As \citet{Baumgartner1993}
and \citet{Jones2004} argue, legislative attention to niche policy
issues operates as a punctuated equilibrium: long periods of equilibrium
(structural silence) are punctuated by episodes of intense attention
driven by exogenous shocks. ESA listing is precisely this type of
niche issue. The 12 active months in our dataset cluster around
identifiable political events --- the 1999 salmon crisis, the 2015 ESA
reauthorisation debates, and the 2025 legislative push --- consistent
with a threshold process rather than a continuous low-probability one.
 
When a model treats structural zeros and sampling zeros as the same type
of observation, the coefficient estimates represent a mixture of two
different relationships: the relationship between news coverage and
structural silence, and the relationship between news coverage and
active count intensity. These are not the same relationship and should
not be estimated as if they were.
 
%% ── 3.7 ─────────────────────────────────────────────────────────
\subsection{The Hurdle model: separating silence from activity}
 
The Hurdle model \citep{Mullahy1986, Zeileis2008} explicitly
represents the two-process structure described above. It consists of
two components estimated simultaneously.
 
\textbf{Part~1} (binary hurdle) models the probability that
$\text{cong}_{t}$ crosses zero --- whether ESA listing crosses the
congressional attention threshold in month~$t$:
 
\begin{equation}
  P(\text{cong}_{t} > 0)
    = \operatorname{logistic}(
        \alpha_0
        + \beta_{40} \cdot \text{news\_lag}
        + \beta_{30} \cdot \text{cong\_lag}
      )
    = \frac{e^{\eta}}{1+e^{\eta}}
  \label{eq:hurdle1}
\end{equation}
 
where $\eta = \alpha_0 + \beta_{40}\cdot\text{news\_lag} +
\beta_{30}\cdot\text{cong\_lag}$. This component uses all 371
observations. The coefficient $\beta_{40}$ directly tests the primary
form of the agenda-setting hypothesis: does media coverage increase the
probability that ESA listing becomes salient enough to elicit
congressional activity?
 
\textbf{Part~2} (count component) models the number of documents
produced given that $\text{cong}_{t} > 0$, estimated from the 12 active
months only:
 
\begin{align}
  \text{cong}_{t} \mid \text{cong}_{t} > 0
    &\;\sim\; \operatorname{Poisson}(\mu_i)
  \label{eq:hurdle2a} \\[3pt]
  \log(\mu_i)
    &= \alpha_1
       + \beta_{41} \cdot \text{news\_lag}_i
       + \beta_{31} \cdot \text{cong\_lag}_i
  \label{eq:hurdle2b}
\end{align}
 
The starting specification uses Poisson in Part~2 rather than Negative
Binomial, following the recommendation of \citet{Zeileis2008}. The
Hurdle model already removes structural zeros through Part~1; adding NB
in Part~2 would introduce a second zero-handling mechanism ($\hat\theta$)
that risks trading off against the hurdle threshold, creating ambiguity
between $\beta_{40}$ and $\beta_{41}$. If residual overdispersion
remains after fitting the Poisson Hurdle, the extension to NB in Part~2
can be evaluated.
 
The full probability mass function combines both components:
 
\begin{equation}
  P(\text{cong}_{t} = k) =
  \begin{cases}
    1 - P(\text{cong}_{t} > 0)
      & k = 0 \\[6pt]
    P(\text{cong}_{t} > 0) \cdot
    \dfrac{P_{\text{Poisson}}(k \mid \mu)}
          {1 - P_{\text{Poisson}}(0 \mid \mu)}
      & k = 1, 2, 3, \ldots
  \end{cases}
  \label{eq:hurdle_pmf}
\end{equation}
 
The two-part structure also produces two distinct coefficients for
news\_lag --- $\beta_{40}$ from Part~1 and $\beta_{41}$ from Part~2 ---
that simultaneously answer two substantively different versions of the
agenda-setting hypothesis:
 
\begin{itemize}[itemsep=2pt]
  \item $\beta_{40}$: Does media coverage predict \textit{whether}
    Congress engages with ESA listing? (the threshold question)
  \item $\beta_{41}$: Does media coverage predict \textit{how much}
    Congress engages, given that it does? (the intensity question)
\end{itemize}
 
No single-component model in this progression --- OLS, Poisson,
Quasi-Poisson, or Negative Binomial --- can answer both questions
simultaneously.
 
The model was fitted in R \citep{RCoreTeam} using the \texttt{pscl}
package \citep{Zeileis2008} with the specification:
 
\begin{verbatim}
hurdle(cong_t ~ news_lag + cong_lag,
       data      = df,
       dist      = "poisson",
       zero.dist = "binomial")
\end{verbatim}
 
%% ── 3.8 ─────────────────────────────────────────────────────────
\subsection{Temporal autocorrelation}
 
Monthly observations are temporally ordered and may not be
independent given only one lag of each variable. Congressional activity
in month $t$ may remain correlated with activity in month $t{+}2$ or
beyond even after conditioning on $\text{cong\_lag}$ at $t{-}1$, and
similarly for the news series \citep{Brandt2000}. The current model
addresses autocorrelation through the inclusion of $\text{cong\_lag}$
as a covariate, following the autoregressive count model approach of
\citet{Brandt2000}. A formal test for residual autocorrelation in the
Hurdle model residuals will be reported in the results. If significant
residual autocorrelation is detected, extensions including a wider lag
window or a correlation structure in the Part~2 component will be
considered in subsequent analyses.
 
%% ════════════════════════════════════════════════════════════════
\section{Discussion}
%% ════════════════════════════════════════════════════════════════
 
%% ── 4.1 ─────────────────────────────────────────────────────────
\subsection{Theoretical implications for agenda-setting}
 
The two-part structure of the Hurdle model maps directly onto the
theoretical framework of punctuated equilibrium theory
\citep{Baumgartner1993, Jones2004}. The logistic component in Part~1
models the punctuation --- the moment at which ESA listing crosses
from equilibrium (congressional silence) into active attention. The
count component in Part~2 models the intensity of activity once
punctuation has occurred. The media coefficient in Part~1 ($\beta_{40}$)
tests whether media coverage is the mechanism that produces punctuation;
the coefficient in Part~2 ($\beta_{41}$) tests whether media coverage
also drives the intensity of activity after punctuation.
 
If both $\beta_{40}$ and $\beta_{41}$ are positive and significant, the
finding would support a strong version of the agenda-setting hypothesis:
media coverage both gets ESA listing onto the congressional agenda and
drives the volume of legislative activity once it is there. If
$\beta_{40}$ is significant but $\beta_{41}$ is not, the finding would
support a threshold version: media matters for attention allocation but
not for production intensity once the issue is salient. These
distinctions cannot be recovered from single-component models.
 
This analysis also speaks to an emerging debate in political
communication about whether agenda-setting operates differently for
niche versus valence issues \citep{Walgrave2008}. ESA listing is
clearly a niche domain: it appears in 17.3\% of news articles in our
corpus but only 5.7\% of congressional documents. For such issues,
the threshold process modelled in Part~1 may be more theoretically
important than the intensity process in Part~2, because the primary
question is whether the issue makes it onto the congressional agenda
at all. The Hurdle model allows both questions to be addressed without
conflating them.
 
%% ── 4.2 ─────────────────────────────────────────────────────────
\subsection{Methodological contributions}
 
This paper contributes a documented model-selection protocol for rare
political event count data that may be useful to researchers facing
similar data structures in other policy domains. The protocol has four
steps. First, characterise the outcome variable before selecting a model:
document the proportion of zeros, the variance-to-mean ratio, the
skewness, and the temporal clustering of non-zero observations. Second,
reject OLS on structural grounds: count data with a non-negative
discrete distribution cannot be appropriately modelled by a continuous
normal regression. Third, evaluate the Poisson family sequentially,
testing for overdispersion at each stage and documenting the specific
assumption that is violated. Fourth, assess whether the zeros are
structural or sampling in nature by examining whether non-zero
observations cluster around identifiable events, and by comparing the
observed zero frequency to the Poisson-predicted frequency at the
estimated mean.
 
This protocol is consistent with the data exploration framework
advocated by \citet{Zuur2010} and the model selection principles in
\citet{Burnham2002}, extended to the specific challenges of political
count data with rare events and structural zeros.
 
A second methodological contribution concerns the integration of
LLM-based narrative coding with statistical inference. The NPF
extraction pipeline \citep{Jones2014} used in this paper represents
a novel approach to operationalising theoretical constructs from
political communication for large-scale quantitative analysis. The
extraction of characters, plots, and morals at scale enables
systematic empirical testing of narrative theory in a way that manual
coding cannot support given practical resource constraints. The
validation procedure described in Section~2 is an important step
toward establishing the reliability of this approach, and future work
should examine the conditions under which LLM-extracted narrative
codes are more or less concordant with human coder judgements.
 
%% ── 4.3 ─────────────────────────────────────────────────────────
\subsection{Limitations and future directions}
 
Several limitations of the current analysis should be acknowledged.
First, the congressional corpus is limited to 300 documents. While
this corpus was stratified to represent the range of ESA policy
activities, the monthly count series for the Species Listing category
contains only 12 active months across 30 years. This sparsity limits
the precision of the count component estimates in Part~2 of the Hurdle
model. Future work should extend the congressional corpus to provide
a more comprehensive monthly count series.
 
Second, the model includes only two predictors: news\_lag and
cong\_lag. Congressional attention to ESA listing is likely influenced
by other factors that are not captured in the current specification,
including election cycles, changes in committee leadership, and
executive branch actions such as proposed listing decisions by the
Fish and Wildlife Service. The inclusion of these covariates in both
parts of the Hurdle model would strengthen causal identification
and address potential confounding.
 
Third, the autocorrelation structure is addressed only through a
single lag. As noted in Section~3.8, residual autocorrelation may
remain in the Hurdle model residuals. Formal testing and, if necessary,
modelling of this autocorrelation structure is an important next step.
 
Fourth, this paper focuses on ESA Species Listing. The other nine
ESA meta-categories --- including Legislative Action, Regulatory Action,
Litigation, and Habitat/Recovery --- each represent distinct policy
processes with potentially different agenda-setting dynamics. A
systematic analysis across all categories would provide a more complete
picture of how media narration of ESA events relates to the congressional
policy process.

% Entries for the entire Anthology, followed by custom entries
\bibliography{anthology,custom}
\bibliographystyle{acl_natbib}

\end{document}
